%------------------------------------------------------------------------------%
\documentclass[fleqn,aspectratio=169,12pt]{beamer}

\usepackage{gaal}

\title{Elipse}

%------------------------------------------------------------------------------%
\begin{document}

\addtitle

%------------------------------------------------------------------------------%
\section{Equação Canônica da Elipse}

%------------------------------------------------------------------------------%
\begin{frame}{Equação Canônica da Elipse}
  Equação \emph{canônica da elipse}
  \[
    \frac{x^2}{a^2} + \frac{y^2}{b^2} = 1
  \]
  com \(a > 0\) e \(b > 0\)

  \bigskip
  Elipse centrada na origem e alinhada com os eixos $x$ e $y$
\end{frame}

%------------------------------------------------------------------------------%
\begin{frame}{Equação Matricial da Elipse}
  Equação \emph{matricial da elipse}
  \[
    v^t A v = 1
  \]
  com
  \[
    v = \begin{bmatrix}
      x \\
      y
    \end{bmatrix}
    \qquad\qquad
    A = \begin{bmatrix}
      \dfrac{1}{a^2} & 0              \\[2mm]
      0              & \dfrac{1}{b^2}
    \end{bmatrix}
    \qquad\qquad
    a > 0
    \qquad
    b > 0
  \]
\end{frame}

%------------------------------------------------------------------------------%
\begin{proofframe}{Pontos Relevantes da Elipse}
\begin{columns}[t]
\column{0.4\textwidth}
  Fazendo \(y=0\)
  \begin{align*}
    \frac{x^2}{a^2} + \frac{y^2}{b^2} & = 1     \\
    \frac{x^2}{a^2} + \frac{0}{b^2}   & = 1     \\
    \frac{x^2}{a^2}                   & = 1     \\
    x^2                               & = a^2   \\
    \abs{x}                           & = a     \\
    x                                 & = \pm a
  \end{align*}

\column{0.4\textwidth}
  Fazendo \(x=0\)
  \begin{align*}
    \frac{x^2}{a^2} + \frac{y^2}{b^2} & = 1     \\
    \frac{0}{a^2} + \frac{y^2}{b^2}   & = 1     \\
    \frac{y^2}{b^2}                   & = 1     \\
    y^2                               & = b^2   \\
    \abs{y}                           & = b     \\
    y                                 & = \pm b
  \end{align*}
\end{columns}
\end{proofframe}

%------------------------------------------------------------------------------%
\begin{frame}{Elementos Geométricos da Elipse}
  \emph{Centro} da elipse
  \[
    C = (0, 0)
  \]

  A elipse cruza os eixos nos pontos
  \[
    (-a, 0) \qquad
    ( a, 0) \qquad
    (0, -b) \qquad
    (0,  b)
  \]

  Se \(a>b\), o eixo maior é horizontal

  Se \(b>a\), o eixo maior é vertical

  Se \(a=b\), a equação representa uma circunferência
\end{frame}

%------------------------------------------------------------------------------%
\begin{frame}{Gráfico da Elipse}{}
{\color{red} Incluir gráficos}
\end{frame}

%------------------------------------------------------------------------------%
%------------------------------------------------------------------------------%
\begin{question}
  Dada a equação da elipse
  \[
    \frac{x^2}{9}+\frac{y^2}{4} = 1
  \]
  \vspace{-1.5\baselineskip}
  \begin{enumerate}
    \item identifique os semieixos ($a$ e $b$)
    \item escreva a equação matricial
    \item esboce a curva
  \end{enumerate}
\end{question}

%------------------------------------------------------------------------------%
\begin{resolution}{Semieixos $a$ e $b$}
  Comparando
  \[
    \frac{x^2}{9} + \frac{y^2}{4} = 1
  \]
  \pause
  com a equação canônica
  \[
    \frac{x^2}{a^2} + \frac{y^2}{b^2} = 1
  \]
  \pause
  identificamos
  \[
    \pause
    a^2 = 9
    \qquad
    \pause
    a = 3
  \]
  \[
    \pause
    b^2 = 4
    \qquad
    \pause
    b = 2
  \]
\end{resolution}

%------------------------------------------------------------------------------%
\begin{resolution}{Equação Matricial}
  Equação Matricial
  \[
    v^t A v = 1
  \]
  \pause
  com
  \[
    A
    = \begin{bmatrix}
      \dfrac{1}{a^2} & 0              \\[2mm]
      0              & \dfrac{1}{b^2}
    \end{bmatrix}
    \pause
    = \begin{bmatrix}
      \dfrac{1}{9} & 0            \\[2mm]
      0            & \dfrac{1}{4}
    \end{bmatrix}
  \]
\end{resolution}

%------------------------------------------------------------------------------%
\begin{resolution}{Esboço}
  \centering
  \includegraphics{I-04-Elipse-ex-1}
\end{resolution}

%------------------------------------------------------------------------------%
\endinput


%------------------------------------------------------------------------------%
\section{Translação da Elipse}

%------------------------------------------------------------------------------%
\begin{frame}{Translação da Elipse}
  Novo centro \(C = (h, k)\)

  Novas coordenadas
  \[
    X = x - h
    \qquad
    Y = y - k
  \]

  Equação canônica nas novas coordenadas
  \[
    \frac{X^2}{a^2} + \frac{Y^2}{b^2} = 1
  \]
  \[
    \frac{(x-h)^2}{a^2} + \frac{(y-k)^2}{b^2} = 1
  \]
\end{frame}

%------------------------------------------------------------------------------%
\begin{proofframe}{Equação Matricial da Elipse Transladada}
  Novo de centro
  \[
    C
    = \begin{bmatrix}
      h \\
      k
    \end{bmatrix}
  \]

  \bigskip
  Novas coordenadas
  \[
    V
    = \begin{bmatrix}
      X \\
      Y
    \end{bmatrix}
    = v - C
    = \begin{bmatrix}
      x \\
      y
    \end{bmatrix}
    -
    \begin{bmatrix}
      h \\
      k
    \end{bmatrix}
    = \begin{bmatrix}
      x - h \\
      y - k
    \end{bmatrix}
  \]
\end{proofframe}

%------------------------------------------------------------------------------%
\begin{proofframe}{Equação Matricial da Elipse Transladada}
  \begin{align*}
    \emph{V^t A V}    & = \emph{1} \\
    (v-C)^t A (v-C)   & = 1 \\
    \begin{bmatrix}
      x - h \\
      y - k
    \end{bmatrix}^t
    \begin{bmatrix}
      \sfrac{1}{a^2} & 0              \\
      0              & \sfrac{1}{b^2}
    \end{bmatrix}
    \begin{bmatrix}
      x - h \\
      y - k
    \end{bmatrix} & = 1 \\
    \emph{\frac{(x-h)^2}{a^2} + \frac{(y-k)^2}{b^2}} & = \emph{1} \\
    \frac{b^2(x-h)^2 + a^2(y-k)^2}{a^2b^2}    & = 1 \\
    b^2(x-h)^2 + a^2(y-k)^2                   & = a^2b^2
  \end{align*}
\end{proofframe}

%------------------------------------------------------------------------------%
\begin{proofframe}{Equação Matricial da Elipse Transladada}
  \begin{align*}
    b^2(x-h)^2 + a^2(y-k)^2                             & = a^2b^2 \\
    b^2(x^2 -2hx + h^2) + a^2(y^2 - 2ky + k^2)          & = a^2b^2 \\
    b^2x^2 - 2b^2hx + b^2h^2 + a^2y^2 - 2a^2ky + a^2k^2 & = a^2b^2 \\
    b^2x^2 + a^2y^2 - 2b^2hx - 2a^2ky + b^2h^2 + a^2k^2 & = a^2b^2 \\
    \begin{bmatrix}
      x \\
      y
    \end{bmatrix}^t
    {\color{magenta}\begin{bmatrix}[c]
      b^2 & 0   \\
      0   & a^2
    \end{bmatrix}}
    \begin{bmatrix}
      x \\
      y
    \end{bmatrix}
    +
    {\color{purple}\begin{bmatrix}
      - 2b^2h \\
      - 2a^2k
    \end{bmatrix}^t}
    \begin{bmatrix}
      x \\
      y
    \end{bmatrix}
    + {\color{violet} (b^2h^2 + a^2k^2 - a^2b^2)} & = 0 \\
    v^t {\color{magenta} M} v
    + {\color{purple} B^t}v
    + {\color{violet} f} & = 0
  \end{align*}
\end{proofframe}

%------------------------------------------------------------------------------%
%------------------------------------------------------------------------------%
\begin{question}
  Dada a equação da elipse
  \[
    \frac{(x-2)^2}{16} + \frac{(y+1)^2}{9} = 1
  \]
  \vspace{-1.5\baselineskip}
  \begin{enumerate}
    \item identifique o centro e os semieixos ($a$ e $b$)
    \item escreva a equação matricial
    \item esboce a curva
  \end{enumerate}
\end{question}

%------------------------------------------------------------------------------%
\begin{resolution}{Semieixos $a$ e $b$}
  Comparando
  \[
    \frac{(x-2)^2}{16} + \frac{(y+1)^2}{9} = 1
  \]
  \pause
  com a equação canônica transladada
  \[
    \frac{(x-h)^2}{a^2} + \frac{(y-k)^2}{b^2} = 1
  \]
  \pause
  identificamos
  \(
    \qquad
    \pause
    a^2 = 16
    \qquad
    \pause
    a = 4
    \qquad\qquad
    \pause
    b^2 = 9
    \qquad
    \pause
    b = 3
  \)

  \pause
  centro:
  \(
    \qquad
    \pause
    C
    = (h, k)
    \pause
    = (2, -1)
  \)
\end{resolution}

%------------------------------------------------------------------------------%
\begin{resolution}{Equação Matricial}
  Equação Matricial
  \[
    V^t A V = 1
  \]
  \pause
  com
  \[
    A
    = \begin{bmatrix}
      \dfrac{1}{a^2} & 0              \\[2mm]
      0              & \dfrac{1}{b^2}
    \end{bmatrix}
    \pause
    = \begin{bmatrix}
      \dfrac{1}{16} & 0            \\[2mm]
      0             & \dfrac{1}{9}
    \end{bmatrix}
  \]
  \[
    \pause
    V
    = \begin{bmatrix}
      x - h  \\
      y - k
    \end{bmatrix}
    \pause
    = \begin{bmatrix}[c]
      x - 2  \\
      y - (-1)
    \end{bmatrix}
    \pause
    = \begin{bmatrix}
      x - 2  \\
      y + 1
    \end{bmatrix}
  \]
\end{resolution}

%------------------------------------------------------------------------------%
\begin{resolution}{Esboço}
  \centering
  \includegraphics{I-04-Elipse-ex-2}
\end{resolution}

%------------------------------------------------------------------------------%
\endinput

%------------------------------------------------------------------------------%
\begin{question}
  Determine o conjunto de soluções reais da equação
  \[
    \frac{(x-1)^2}{4} + \frac{(y+2)^2}{9} = 0
  \]

\end{question}

%------------------------------------------------------------------------------%
\begin{resolution}{Solução}
  Cada parcela é não negativa
  \[
    \frac{(x-1)^2}{4}\geq 0
    \qquad
    \frac{(y+2)^2}{9}\geq 0
  \]
  A soma é nula somente se
  \[
    x - 1 = 0
    \qquad
    y + 2 = 0
  \]
  Logo o conjunto de soluções é
  \[
    S = \{(1,-2)\}
  \]
\end{resolution}

%------------------------------------------------------------------------------%
\endinput


%%------------------------------------------------------------------------------%
%\section{Lista Mínima}
%
%%------------------------------------------------------------------------------%
%\begin{frame}{Lista Mínima}
%  \begin{enumerate}
%    \item
%  \end{enumerate}
%
%  \vfill
%  Atenção: A prova é baseada no livro, não nas apresentações
%\end{frame}

%------------------------------------------------------------------------------%
\end{document}
%------------------------------------------------------------------------------%
